\documentclass[11pt]{article}
\usepackage[margin=1in]{geometry}
\usepackage{amsmath,amssymb}
\usepackage{booktabs}
\usepackage{enumitem}
\usepackage[T1]{fontenc}
\usepackage{lmodern}
\usepackage{microtype}
\usepackage[hidelinks]{hyperref}
\usepackage{xcolor}

\hypersetup{
  pdftitle={Modeling Xenon Delivery by Saturated-Liquid Replacement},
  pdfauthor={},
  pdfsubject={Mass transfer, liquid replacement, and spherical tissue diffusion}
}
\setlength{\parskip}{0.55em}
\setlength{\parindent}{0pt}
\setlist{nosep}

\title{\textbf{Modeling Xenon Delivery by Saturated-Liquid Replacement}\\
\large A mass-transfer framework for comparison with a sealed headspace}
\date{4 November 2026}

\begin{document}
\maketitle

\begin{abstract}
This note formulates a xenon-delivery experiment in which liquid is
equilibrated with xenon in an external saturator and then flows into a vessel
containing an organoid. The vessel liquid is progressively replaced rather
than pressurized by a sealed headspace. Two vessel descriptions are given:
a well-mixed control volume and a spatial advection--dispersion model. Both
can be coupled to radial xenon diffusion in a spherical organoid. The
assumptions, boundary conditions, equilibrium predictions, and measurements
needed for comparison are stated explicitly.
\end{abstract}

\section{Physical picture and modeling distinction}

The incoming and resident liquids are generally the \emph{same solvent};
they are not two immiscible fluids with a sharp material interface. The
quantity being transported through the vessel is dissolved xenon. Bulk liquid
motion carries xenon by advection, while molecular diffusion and
flow-induced mixing smooth concentration gradients. A sharp replacement
front may be a useful approximation at high P\'eclet number, but is not
required by the model.

The external saturator supplies a prescribed xenon concentration. Liquid
replacement changes the vessel concentration, and xenon then diffuses from
the liquid into the organoid. Unlike a sealed-headspace experiment, the
incoming liquid is an ongoing source: the vessel's final concentration is
not found by conserving the initial sealed-vessel inventory.

\section{Xenon concentration in the incoming liquid}

Let \(P_s(t)\) be the xenon partial pressure in the saturator, \(T\) the
absolute temperature, and \(R_g\) the ideal-gas constant. At gas--liquid
equilibrium, Henry's law in concentration-over-pressure convention gives
\begin{equation}
  C_{\mathrm{in}}(t) = H_{cp}(T)\,P_s(t).
  \label{eq:henry}
\end{equation}
Here \(H_{cp}\) has units of concentration per pressure. If using the
project's partition notation, the same relation may be written
\begin{equation}
  C_{\mathrm{in}}(t) =
  \lambda_l \frac{P_s(t)}{R_gT},
  \label{eq:project-henry}
\end{equation}
provided \(\lambda_l\) is defined consistently with that convention.
Henry-constant conventions differ (some use pressure over mole fraction,
others concentration over pressure); units and definition must therefore
accompany any numerical value \cite{IUPACHenry,Sander2023}.

If saturation is not complete, equation~\eqref{eq:henry} is an equilibrium
target rather than the actual outlet concentration. A lumped, well-mixed
saturator can be approximated by
\begin{equation}
  \frac{dC_s}{dt}
  = k_La\left(H_{cp}P_s-C_s\right),
  \qquad C_{\mathrm{in}}=C_s,
  \label{eq:saturator}
\end{equation}
where \(C_s\) is the liquid concentration in the saturator, \(k_La\) is a
volumetric gas--liquid transfer coefficient, and \(P_s\) is assumed known.
This first-order equation is a reduced model, not a universal saturator
equation; an actual contactor may instead need a plug-flow or staged
gas--liquid model.

\section{Vessel option A: well-mixed liquid}

If mixing in the vessel is fast compared with both replacement and uptake,
represent the vessel by one concentration \(C_v(t)\). For constant liquid
volume \(V_l\), inflow and outflow \(Q\), and one spherical organoid of
radius \(R\), its xenon balance is
\begin{equation}
  V_l\frac{dC_v}{dt}
  = Q\left(C_{\mathrm{in}}-C_v\right)
    - A_o J(t),
  \qquad A_o=4\pi R^2.
  \label{eq:cstr}
\end{equation}
The first term on the right is xenon delivered by flow minus xenon removed
by outflow. The second is uptake by the organoid. \(J>0\) denotes xenon flux
from liquid into tissue, with units \(\mathrm{mol\,m^{-2}\,s^{-1}}\).
Equation~\eqref{eq:cstr} assumes equal inflow and outflow so that \(V_l\)
stays constant.

If there is no tissue uptake, the response to a constant inlet concentration
is
\begin{equation}
  C_v(t)=C_{\mathrm{in}}+
  \left[C_v(0)-C_{\mathrm{in}}\right]e^{-t/\tau_{\mathrm{flow}}},
  \qquad \tau_{\mathrm{flow}}=\frac{V_l}{Q}.
  \label{eq:replacement}
\end{equation}
Thus \(V_l/Q\) is the ideal well-mixed replacement time. Uptake couples
equation~\eqref{eq:cstr} to the tissue and can delay the vessel concentration
from approaching the inlet value.

\subsection{Tissue diffusion and liquid--tissue coupling}

For a homogeneous spherical organoid with constant effective xenon
diffusivity \(D_t\), let \(C_t(r,t)\) denote tissue concentration. Radial
diffusion obeys
\begin{equation}
  \frac{\partial C_t}{\partial t}
  = \frac{D_t}{r^2}\frac{\partial}{\partial r}
    \left(r^2\frac{\partial C_t}{\partial r}\right),
  \qquad 0<r<R,
  \label{eq:tissue}
\end{equation}
with center symmetry
\begin{equation}
  \left.\frac{\partial C_t}{\partial r}\right|_{r=0}=0.
  \label{eq:center}
\end{equation}

For a perfectly mixed liquid and negligible liquid-side mass-transfer
resistance, assume instantaneous partition equilibrium at the tissue surface:
\begin{equation}
  C_t(R,t)=\lambda_{t/l}C_v(t).
  \label{eq:partition-boundary}
\end{equation}
The corresponding uptake flux is
\begin{equation}
  J(t)=D_t\left.\frac{\partial C_t}{\partial r}\right|_{r=R}.
  \label{eq:tissue-flux}
\end{equation}
With this sign convention, a tissue concentration that is higher at the
surface than in the center has \(\partial_r C_t(R,t)>0\) and takes up xenon.
Equations~\eqref{eq:cstr}--\eqref{eq:tissue-flux} conserve xenon between the
well-mixed vessel and the tissue.

If transport across the liquid boundary layer is not negligible, resolve a
surface liquid concentration \(C_{l,s}\) and use a finite liquid-film
coefficient \(k_l\):
\begin{align}
  J &= k_l\left(C_v-C_{l,s}\right), \label{eq:liquid-film}\\
  C_t(R,t) &= \lambda_{t/l}C_{l,s}(t), \label{eq:finite-partition}\\
  J &= D_t\left.\frac{\partial C_t}{\partial r}\right|_{r=R}.
  \label{eq:finite-flux}
\end{align}
Here \(\lambda_{t/l}\) is the equilibrium tissue-to-liquid concentration
ratio. A non-equilibrium membrane or interfacial resistance can be added as
another resistance in series. These alternatives should not be combined
without checking the definitions of the partition coefficient and transfer
coefficients.

\section{Vessel option B: spatial advection--dispersion}

When gradients along the vessel are important, a one-dimensional
cross-section-averaged model for dissolved xenon is
\begin{equation}
  \frac{\partial C_l}{\partial t}
  + u\frac{\partial C_l}{\partial z}
  = D_{\mathrm{ax}}\frac{\partial^2 C_l}{\partial z^2},
  \qquad 0<z<L_v,
  \label{eq:dispersion}
\end{equation}
where \(u\) is mean liquid velocity and \(D_{\mathrm{ax}}\) is an effective
axial dispersion coefficient. It includes molecular diffusion and, when
appropriate, spreading caused by the velocity profile and unresolved
mixing. This is the axial-dispersion framework; Taylor's analysis explains
how shear flow in a tube can enhance the effective longitudinal spreading
\cite{Taylor1953}.

For a dispersive inlet, Danckwerts boundary conditions are commonly used:
\begin{align}
  uC_l(0,t)-D_{\mathrm{ax}}
  \left.\frac{\partial C_l}{\partial z}\right|_{z=0}
  &=uC_{\mathrm{in}}(t), \label{eq:inlet}\\
  \left.\frac{\partial C_l}{\partial z}\right|_{z=L_v}
  &=0. \label{eq:outlet}
\end{align}
An initial profile \(C_l(z,0)\) must also be specified. These conditions
represent imposed inlet flux and negligible dispersive flux at the outlet.

Equation~\eqref{eq:dispersion} alone describes xenon transport through the
liquid, not uptake by an isolated organoid. For one organoid, the rigorous
spatial model resolves the surrounding liquid and couples its surface flux
to equation~\eqref{eq:tissue}. If many similar organoids are distributed
throughout the vessel and can be represented as a uniform population, a
homogenized sink term may instead be added:
\begin{equation}
  \frac{\partial C_l}{\partial t}
  +u\frac{\partial C_l}{\partial z}
  =D_{\mathrm{ax}}\frac{\partial^2 C_l}{\partial z^2}
   -a_sJ.
  \label{eq:distributed-sink}
\end{equation}
Here \(a_s\) is organoid surface area per liquid volume, and \(J\) is the
local liquid-to-tissue uptake flux. The term \(a_sJ\) has units of
concentration per time. It is not appropriate to spread a single organoid's
uptake uniformly along the vessel unless that is an intentional
volume-averaged approximation.

\section{Equilibrium and step changes}

For a constant, fully equilibrated inlet maintained indefinitely, with no
reaction or degradation, the final state is spatially uniform:
\begin{equation}
  C_v \longrightarrow C_{\mathrm{in}}, \qquad
  C_t(r) \longrightarrow \lambda_{t/l}C_{\mathrm{in}}.
  \label{eq:steady}
\end{equation}
This is the long-time limit of a continuously replenished system. It is
different from the equilibrium of a sealed headspace, where a finite initial
xenon inventory is redistributed among gas, liquid, and tissue. The
equilibrium relation used in the sealed-vessel model should therefore not be
used to set the final concentration of this flow-through experiment.

For equal increments \(\Delta P_s\) in saturator xenon partial pressure, the
ideal inlet concentration after \(k\) increments is
\begin{equation}
  C_{\mathrm{in},k}=H_{cp}\,k\Delta P_s.
  \label{eq:steps}
\end{equation}
If each new level is held indefinitely before the next, its corresponding
liquid and tissue equilibrium levels follow equation~\eqref{eq:steady}.
During finite intervals the actual liquid and tissue concentrations lag
these reference levels. Pressure steps applied to the saturator are not
necessarily instantaneous concentration steps at the vessel inlet if the
saturator or connecting tubing has appreciable dynamics or dead volume.

\section{Dimensionless scales and comparison}

Useful characteristic times and groups include
\begin{align}
  \tau_{\mathrm{flow}} &= \frac{V_l}{Q},
  &\tau_{\mathrm{tissue}} &= \frac{R^2}{D_t},\\
  \mathrm{Pe} &= \frac{uL_v}{D_{\mathrm{ax}}},
  &\mathrm{Fo}_t &= \frac{D_t t}{R^2}.
\end{align}
The flow time estimates replacement in a mixed-volume approximation; the
tissue diffusion time estimates internal equilibration. The P\'eclet number
compares axial advection with effective axial dispersion. A direct comparison
with the sealed-headspace model should use the same organoid parameters,
temperature, initial tissue state, and final inlet xenon level. Then compare
the vessel/interface concentration histories and tissue-average or
center-to-surface concentration histories, rather than comparing nominal
pressures alone.

\section{Recommended first implementation and measurements}

Start with equations~\eqref{eq:cstr}--\eqref{eq:tissue-flux} if the vessel is
well mixed; use equations~\eqref{eq:dispersion}--\eqref{eq:outlet} if
replacement fronts or axial gradients are important. To parameterize either
version, measure or specify:
\begin{enumerate}
  \item \(Q\), \(V_l\), vessel geometry, and the initial liquid xenon level;
  \item saturator pressure, temperature, residence/contact time, and
        preferably inlet xenon concentration (to test complete saturation);
  \item organoid radius, number, exposed area, and placement in the flow;
  \item xenon liquid diffusivity, effective tissue diffusivity, and the
        tissue/liquid partition ratio at the experiment temperature;
  \item whether liquid-side or membrane mass-transfer resistance is
        significant; and
  \item vessel concentration versus time, flow rate, and the timing of each
        saturator pressure change.
\end{enumerate}
Water Henry constants from the literature should not be substituted
automatically for the project's fitted \(\lambda_l\), nor should a tissue
partition factor be assumed universal. Both are temperature-, medium-, and
definition-dependent quantities \cite{Sander2023,Ercan1979}.

\section*{Notation}
\begin{center}
\begin{tabular}{@{}ll@{}}
\toprule
Symbol & Meaning\\
\midrule
\(C_{\mathrm{in}},C_v,C_l,C_t\) & inlet, vessel, liquid-field, tissue xenon concentrations\\
\(P_s\) & xenon partial pressure in the external saturator\\
\(H_{cp}\) & Henry solubility coefficient, concentration per pressure\\
\(\lambda_{t/l}\) & tissue-to-liquid equilibrium concentration ratio\\
\(Q,V_l\) & liquid volumetric flow rate and vessel liquid volume\\
\(D_{\mathrm{ax}},D_t\) & effective axial dispersion and tissue diffusivity\\
\(J\) & xenon uptake flux from liquid into tissue\\
\(R,L_v,u\) & organoid radius, vessel length, and mean liquid velocity\\
\bottomrule
\end{tabular}
\end{center}

\begin{thebibliography}{9}
\bibitem{ProjectPDF}
G.~L. Timmerman,
\emph{Xenon diffusion in fluid},
project manuscript supplied as \texttt{Xenon\_diffusion.pdf}, April 2026.

\bibitem{IUPACHenry}
International Union of Pure and Applied Chemistry (IUPAC),
``Henry's law,'' \emph{Compendium of Chemical Terminology (the Gold Book)},
\href{https://goldbook.iupac.org/terms/view/R05135}
     {https://goldbook.iupac.org/terms/view/R05135}.

\bibitem{Sander2023}
R. Sander,
``Compilation of Henry's law constants (version 5.0.0) for water as solvent,''
\emph{Atmospheric Chemistry and Physics}, vol.~23, pp.~10901--12440, 2023.
\href{https://doi.org/10.5194/acp-23-10901-2023}
     {doi:10.5194/acp-23-10901-2023}.

\bibitem{Taylor1953}
G.~I. Taylor,
``Dispersion of soluble matter in solvent flowing slowly through a tube,''
\emph{Proceedings of the Royal Society of London. Series A},
vol.~219, no.~1137, pp.~186--203, 1953.
\href{https://doi.org/10.1098/rspa.1953.0139}
     {doi:10.1098/rspa.1953.0139}.

\bibitem{Crank1975}
J. Crank,
\emph{The Mathematics of Diffusion}, 2nd ed.
Oxford: Clarendon Press, 1975.

\bibitem{Ercan1979}
M.~T. Ercan,
``Solubility coefficients of \({}^{133}\mathrm{Xe}\) in water, saline, dog
blood and organs,'' \emph{International Journal of Applied Radiation and
Isotopes}, vol.~30, pp.~757--759, 1979.
\href{https://doi.org/10.1016/0020-708X(79)90156-X}
     {doi:10.1016/0020-708X(79)90156-X}.
\end{thebibliography}

\end{document}
